% sec-04-trial-protocol.tex - Section 4, Trial Protocol.  A main section.
% DRAFT stage.  Four figure slots, Figures 10 to 13, and three tables, Tables
% 11, 12 and 13.  Sources are trial-protocol/final-protocol/publication and
% trial-phase-2/final-protocol/publication/author.
\section{Trial Protocol}
\label{sec:protocol}

\subsection{The Phase 1 study}

The Phase 1 protocol, version 1.0.0, deposited June 21, 2026, is an open-label,
single-arm, dose-finding study using a standard 3+3 escalation across three
daraxonrasib dose levels with staggered sentinel enrollment, planned for up to
n = 18 participants \cite{onpremwhippl}. A mandatory Phase 0 simulation
validation gate precedes any patient-facing robotic activity: at least 1000
simulated procedures across at least two independent frameworks, with a Unified
Safety Level rating at or above 7.0. Table 11 is the deposited synopsis in
full.

\draftinstr{Stage 7: quote directly from
\appfile{trial-protocol/final-protocol/publication/LaTeX Source Files.zip},
sections/sec-01-summary.tex, the Design paragraph and the Intervention
paragraph, including the 10 kHz heartbeat broadcast bus, the cross-arm
emergency-stop budget at or below 3 ms, the per-arm tip-force cap at or below
3 N, the cumulative cross-arm cap at or below 18 N, and vascular no-fly gating
around the superior mesenteric and portal veins.}

\begin{apptable}
\begin{tabularx}{\textwidth}{>{\raggedright\arraybackslash}p{3.2cm} >{\raggedright\arraybackslash}X}
\headrow \thead{Element} & \thead{Phase 1, first-in-human, v1.0.0, June 21, 2026}\\
\midrule
Design & Open-label, single-arm, 3+3 escalation, DL1 to DL3, staggered sentinel enrollment \\
Sample size & Up to n = 18 \\
Primary endpoints & DLT rate during escalation; device- and procedure-related SAE rate through day 30, including Clavien-Dindo grade III or higher \\
Key secondary & R0 rate at day-90 pathology; ISGPS grade B or C fistula rate; 90-day mortality; PFS; OS to 24 months \\
Population & Adults, resectable or borderline-resectable KRAS G12-mutated PDAC, ECOG 0 or 1 \\
Gate & Phase 0: at least 1000 simulations, at least 2 frameworks, USL at or above 7.0 \\
Regulatory basis & IND 21 CFR \S 312 plus IDE 21 CFR \S 812, Physical AI Subpart J \\
\end{tabularx}
\end{apptable}
\vspace{-0.6cm}
\tabcap{Table 11. The Phase 1 protocol synopsis as deposited, with the gate\\
that must close before any patient-facing robotic activity begins.}

\figslot{plantuml-type / state with guards}{9.6}{Figure 10. One participant from screening to the 24-month endpoint,
with the\\guard on every transition and the two conditions that route to a study
hold.}

Figure 10 takes the participant's point of view. Nine forward states run down a
left column at a constant pitch, two exception states sit in a right column, and
every one of the fourteen transitions carries a guard written as a quantity that
can be evaluated at the bedside. The guards out of any one state partition the
space, so no participant can occupy two successors at once and no state can
deadlock.

\subsection{The Phase 2 study}

\draftinstr{Stage 7: quote directly from
\appfile{trial-phase-2/final-protocol/publication/author/LaTeX Source Files.zip},
sections/sec-01-summary.tex, the Design paragraph: multicenter, randomized 1:1,
parallel-group, controlled, open-label, BICR of imaging, central masked
pathology, blinded endpoint adjudication, eight high-volume HPB centers, n = 220,
stratified by resectability, KRAS allele, neoadjuvant therapy and site;
approximately 140 PFS events for 85 percent power at two-sided alpha 0.05 to
detect a hazard ratio of 0.60, median PFS 14.0 months control against 23.3
months experimental, one interim analysis at about 60 percent of events using
O'Brien-Fleming spending.}

\begin{apptable}
\begin{tabularx}{\textwidth}{>{\raggedright\arraybackslash}p{3.2cm} >{\raggedright\arraybackslash}X}
\headrow \thead{Element} & \thead{Phase 2, randomized, v1.1.0, June 23, 2026}\\
\midrule
Design & Multicenter, randomized 1:1, parallel-group, controlled, open-label with BICR \\
Sample size & n = 220, approximately 110 per arm, eight high-volume HPB centers \\
Primary endpoint & PFS per RECIST 1.1, investigator assessed, BICR as sensitivity analysis \\
Key secondary & OS; R0 rate; ISGPS grade B or C fistula rate; major pathologic response; KRAS ctDNA clearance at week 12 \\
Statistics & About 140 PFS events, 85 percent power, two-sided alpha 0.05, target hazard ratio 0.60 \\
Stratification & Resectability, KRAS allele, neoadjuvant therapy, site \\
Predicate & Phase 1 v1.0.0, which established the RP2D and the safety architecture \\
\end{tabularx}
\end{apptable}
\vspace{-0.6cm}
\tabcap{Table 12. The Phase 2 protocol synopsis as deposited, and the Phase\\
1 predicate whose result each randomized element of it depends upon.}

Table 12 is the Phase 2 synopsis, and every randomized element in it depends on
a Phase 1 result.

\figslot{mermaid-type / flowchart LR}{7.8}{Figure 11. From the Phase 0 simulation gate through three 3+3 dose
levels to\\Phase 2 randomization, with the quantity that must hold on each rung
shown.}

\figslot{d2-type / containers}{8.4}{Figure 12. The twelve Phase 1 protocol sections sorted into what
Phase 2\\carried unchanged, what it replaced, and what randomization made
necessary.}

Figure 11 draws the clinical progression and Figure 12 draws the documentary
one. Four of the twelve Phase 1 sections were carried into Phase 2 unchanged,
four were replaced outright because the study is no longer single-arm, and four
elements exist in the Phase 2 document only because randomization requires them.
Two days separate the two deposits.

\draftinstr{Stage 7: build Table 13 with columns Document, Version, Deposit
date, Sections, Figures, Elapsed production, Prior-system counterpart. Both rows
plus a combined row. The two-day gap between June 21 and June 23 is the
quantitative point.}

\begin{apptable}
\begin{tabularx}{\textwidth}{>{\raggedright\arraybackslash}p{3.0cm} >{\raggedright\arraybackslash}p{1.5cm} >{\raggedright\arraybackslash}p{2.0cm} >{\raggedright\arraybackslash}p{1.4cm} >{\raggedright\arraybackslash}X}
\headrow \thead{Document} & \thead{Version} & \thead{Deposit} & \thead{Sections} & \thead{Elapsed production against the prior system}\\
\midrule
Phase 1 protocol & v1.0.0 & Jun 21, 2026 & 13 & 2 days against 6 to 12 months \\
Phase 2 protocol & v1.1.0 & Jun 23, 2026 & 13 & 2 days against 6 to 12 months \\
Both, combined & Two & Jun 21 to 23 & 26 & 4 days against a full trial-design cycle \\
\end{tabularx}
\end{apptable}
\vspace{-0.6cm}
\tabcap{Table 13. Both protocol deposits, their versions, character counts and\\figure
counts, and the elapsed production time against the prior system.}

\subsection{The site the protocol assumes}

\figslot{diagrams (python)-type / clustered infrastructure}{9.8}{Figure 13. The four-layer on-premises site stack, and the single
boundary\\across which no inference request, model weight, or patient record
passes.}

Figure 13 draws the equipment a participant is operated within, and the one
boundary that makes the protocol's on-premises claim checkable. The prohibited
path is the only struck-through edge in this paper: a site that cannot show that
edge struck is not running this protocol.

\draftinstr{Stage 7: close the section by connecting Figure 13's boundary to the
consent language in \appfile{trial-protocol/final-protocol/publication},
sections/sec-05-population.tex, and to the autonomy disclosure duty in \S5. This
is a main section: match the character count of sections 3, 5, 6 and 7.}
